\documentclass{subfiles}
\begin{document}
    In the previous section we showed that red-black trees are good search trees,
        since the tree height is proportional to the black height of the root. 
        In the present section we focus on another problem: how to mantain balanced a RB tree?    
        To do so, we introduce the notion of rotation.

    When inserting or deleting a node from a red-black tree,
        it can happen that the properties of RB trees are no loger satisfied;
        we shall see in fact, that adding a new node to the tree breaks either property 
        \ref{RB-trees:const2}  or property \ref{RB-trees:const4}.
        Hence in addition to the procedure to add a node, 
        we need one that restores all the properties. 
        Such procedure, which we refer to as \lstinline{RB-Insert-Fixup},
        make use of rotations which are local operation on pointers. 
        
        \subfile{../../extras/TikZ/Figure - Tree rotations}

    In \Cref{Fig:1} we show that there exist two fundamental kinds of rotations: 
        to the left and to the right. Let us also note that it can be proved that at most 
        two rotations are needed per operation. We briefly overview \lstinline{RotateLeft},
        as this is easely understood and \lstinline{RotateRight} is symmetrical.

    Recall that the elements of the red-black tree satisfy some order relation;
        that is, all elements on the left side of a node have a key that is lower 
        or equal to that of the node itself. In what follows we use the following 
        notation: \(x.key\) represents \(x\)'s key, \(x.left\) (respectively \(x.right\))
        to refer to the left (the right) child, \(x.p\) its parent and \(x.p.p\)
        is grandparent.

    Focusing on \Cref{Fig:1}, we have: that \(x.key < y.key\) and \(\beta.key > x.key\).
        Therefore, if we make the subtree rooted at \(\beta\) a right subtree of \(x\)
        the order is still satisfied. The only issue is that both \(y\) and \(x\)
        have \(\beta\) as child, loosing the tree structure. What we need is simply to 
        make \(x.p = y\) and \(y.left = x\).


    \paragraph*{RB trees: node insertion}
    \subfile{../paragraph/2.1.2.a - RB tree insertion}
   
    \paragraph*{RB trees: node deletion}
    \subfile{../paragraph/2.1.2.b - RB tree deletion.tex}
\end{document}
